\subsection{测度分解为具有紧支集的测度之和}

命题 4. --- 设 \(\mu\) 是局部紧空间 \(T\) 上的一个正测度，并设 \(\mathfrak{K}\) 是紧子集的一个 \(\mu\) -稠密集（\(T\) 的）。存在正测度的一个可和族 \((\mu_{\alpha})_{\alpha \in A}\) （\(T\) 上的），使 \(\mu = \sum_{\alpha \in A} \mu_{\alpha}\) ，且使诸测度 \(\mu_{\alpha}\) 的支集属于 \(\mathfrak{K}\) 并构成两两不相交紧集的一个局部可数族。

若测度 \(\mu\) 是适度的，则可取指标集 A 为可数的。

考虑 K 的两两不相交元素的一个局部可数族 \((\mathrm{K}_{\alpha})_{\alpha\in\mathrm{A}}\) ，使集 \(N = T - \bigcup K_{\alpha}\) 局部 \(\mu\) -可忽略（第四章，§5 第 9 号命题 14）。对每个函数 \(f \in \mathcal{K}(\mathrm{T})\) ，令

\[
\mu_ {\alpha} (f) = \mu (f \varphi_ {\mathrm{K} _ {\alpha}});
\]

线性形式 \(\mu_{\alpha}\) （\(\mathcal{K}(\mathrm{T})\) 上的）是正的，因而是一个正测度，其支集含于 \(\mathrm{K}_{\alpha}\) 。由于含于 \(\mathfrak{K}\) 的一个元素之中的每个紧集都属于 \(\mathfrak{K}\) ，故 \(\operatorname{Supp}(\mu_{\alpha}) \in \mathfrak{K}\) 对一切 \(\alpha \in \mathrm{A}\) 成立。剩下只需证明族 \((\mu_{\alpha})\) 可和且其和等于 \(\mu\) ，换言之，\(\sum_{\alpha \in \mathrm{A}} \mu_{\alpha}(f) = \mu(f)\) 对每个函数 \(f \in \mathcal{K}_{+}(\mathrm{T})\) 成立。设 S 是 \(f\) 的（紧）支集，并设 \(\mathrm{A}'\) 是由 \(\alpha \in \mathrm{A}\) （使 \(\mathrm{S} \cap \mathrm{K}_{\alpha} \neq \varnothing\) 者）组成的可数集。由于集 \(\mathrm{N} \cap \mathrm{S}\) 是 \(\mu\) -可忽略的，

\[
\begin{array}{c} \mu (f) = \mu (f \varphi_ {\mathrm{S}}) = \sum_ {\alpha \in \mathrm{A} ^ {\prime}} \mu (f \varphi_ {\mathrm{S} \cap \mathrm{K} _ {\alpha}}) = \sum_ {\alpha \in \mathrm{A} ^ {\prime}} \mu (f \varphi_ {\mathrm{K} _ {\alpha}}) \\ = \sum_ {\alpha \in \mathrm{A}} \mu (f \varphi_ {\mathrm{K} _ {\alpha}}) = \sum_ {\alpha \in \mathrm{A}} \mu_ {\alpha} (f). \end{array}
\]

这就完成了一般情形的证明。若 \(\mu\) 是适度的，则集 T 是 \(\mu\) -适度的，于是 T 是紧集的一个序列 \((\mathrm{L}_n)\) 与一个可忽略集之并（§1 第 2 号命题 5）；设 \(\mathbf{A}'\) 是由 \(\alpha \in \mathbf{A}\) （使 \(\mathrm{K}_{\alpha}\) 与某个 \(\mathrm{L}_n\) 相交者）组成的可数集。则 \(\mu_{\alpha} = 0\) （对 \(\alpha \notin \mathbf{A}'\) ），而命题的最后一句话立即得证。

注记. --- 一个正测度可能是一个具有紧支集的测度序列之和，而不是适度的（参看 §1 习题 4 a)）。

